\section{Number Theory}
\subsection{GCD Properties}
In particular, recalling that gcd is a positive integer valued function we obtain that $\gcd(a, b \cdot c) = 1$ if and only if $\gcd(a, b) = 1$ and $\gcd(a, c) = 1$.

15. $\gcd(a_1, \ldots, a_r) = \gcd(\gcd(a_1, \ldots, a_{r-1}), a_r)$ when none $=0$

16. $\gcd(a,b) \cdot \text{lcm}(a,b) = |a \cdot b|$, so for any two numbers LCM is always multiple of GCD.

17. In a Cartesian coordinate system, $\gcd(a, b)$ can be interpreted as the number of segments between points with integral coordinates on the straight line segment joining the points $(0, 0)$ and $(a, b)$.

18. Two numbers are called relatively prime, or coprime, if their greatest common divisor equals 1. For example, 9 and 28 are relatively prime.

19. Any two consecutive integers will be Co-Prime. Ex: 25,26 is co-prime. 5,6 is co-prime and so on.

20. Two or more odd consecutive positive integer GCD is also 1.
\begin{minted}{c++}
int gcdFast(int a, int b) {
  int result;
  __asm__ __volatile__(
    "movl %1, %%eax;"
    "movl %2, %%ebx;"
    "CONTD: cmpl $0, %%ebx;"
    "je DONE;"
    "xorl %%edx, %%edx;"
    "idivl %%ebx;"
    "movl %%ebx, %%eax;"
    "movl %%edx, %%ebx;"
    "jmp CONTD;"
    "DONE: movl %%eax, %0;"
    : "=g"(result)
    : "g"(a), "g"(b));
  return result;
}
\end{minted}
\subsection{Extended Euclid and CRT}
\begin{minted}{c++}
using T = __int128_t;
T egcd(T a, T b, T &x, T &y) {
  T xx = y = 0, yy = x = 1;
  while (b) {
    T q = a / b;
    T t = b; b = a % b; a = t;
    t = xx; xx = x - q * xx; x = t;
    t = yy; yy = y - q * yy; y = t;
  }
  return a;
}
pair<T, T> crt(T a1, T m1, T a2, T m2) {
  T p, q;
  T g = egcd(m1, m2, p, q);
  if (a1 % g != a2 % g) return {0, -1};
  T m = m1 / g * m2;
  p = (p % m + m) % m;
  q = (q % m + m) % m;
  T x = p * a2 % m * (m1 / g) % m +
        q * a1 % m * (m2 / g) % m;
  return {x % m, m};
}
\end{minted}
\subsection{Linear Diophantine Equation}
Solves $Ax + By = C$. Solutions exist iff
$g = \gcd(|A|, |B|)$ divides $C$.
General: $x = x_0 + k \frac{B}{g}$, $y = y_0 - k \frac{A}{g}$.
\begin{minted}{c++}
ll extgcd(ll a, ll b, ll &x, ll &y) {
  if (b == 0) { x = 1; y = 0; return a; }
  ll x1, y1;
  ll d = extgcd(b, a % b, x1, y1);
  x = y1; y = x1 - y1 * (a / b);
  return d;
}
bool find_any_solution(ll a, ll b, ll c,
ll &x0, ll &y0, ll &g) {
  g = extgcd(abs(a), abs(b), x0, y0);
  if (c % g) return false;
  x0 *= c / g; y0 *= c / g;
  if (a < 0) x0 = -x0;
  if (b < 0) y0 = -y0;
  return true;
}
\end{minted}
\subsection{SPF, Möbius, Primes in O(N)}
\begin{minted}{c++}
const int N = 1e6 + 2;
vector<int> spf(N), mob(N), primes;
void build() {
  mob[1] = 1;
  for (int i = 2; i < N; i++) {
    if (spf[i] == 0) {
      spf[i] = i, mob[i] = -1,
      primes.push_back(i);
    }
    for (int j = 0; i * primes[j] < N; j++) {
      spf[i * primes[j]] = primes[j];
      if (i % primes[j])
        mob[i * primes[j]] = mob[i] * mob[primes[j]];
      else mob[i * primes[j]] = 0;
      if (primes[j] == spf[i]) break;
    }
  }
}
\end{minted}
\subsection{Segmented Sieve}
\begin{minted}{c++}
vector<char> segmentedSieve(long long L,
long long R) {
  long long lim = sqrt(R);
  vector<char> mark(lim + 1, false);
  vector<long long> primes;
  for (long long i = 2; i <= lim; ++i) {
    if (!mark[i]) {
      primes.emplace_back(i);
      for (long long j = i * i; j <= lim; j += i)
        mark[j] = true;
    }
  }
  vector<char> isPrime(R - L + 1, true);
  for (long long i : primes) {
    ll sm = (L / i) * i;
    if (sm < L) sm += i;
    for (long long j = max(i * i, sm); j <= R; j += i)
      isPrime[j - L] = false;
  }
  if (L == 1) isPrime[0] = false;
  return isPrime;
}
\end{minted}
\subsection{Phi (Euler Totient)}
\begin{minted}{c++}
int phi(int n) {
  int result = n;
  for (int i = 2; i * i <= n; i++) {
    if (n % i == 0) {
      while (n % i == 0) n /= i;
      result -= result / i;
    }
  }
  if (n > 1) result -= result / n;
  return result;
}
void phi_1_to_n() {
  phi[1] = 1;
  for (int i = 2; i < N; i++) {
    if (phi[i] == 0) {
      primes.push_back(i);
      phi[i] = i - 1;
    }
    for (int p : primes) {
      if (i * p >= N) break;
      if (i % p == 0) {
        phi[i * p] = phi[i] * p; break;
      } else {
        phi[i * p] = phi[i] * (p - 1);
      }
    }
  }
}
\end{minted}
\subsection{Miller-Rabin Primality Test}
\begin{minted}{c++}
ull pow(ull a, ull t, ull mod) {
  ull r = 1;
  for (; t; t >>= 1, a = (ull)a * a % mod)
    if (t & 1) r = (ull)r * a % mod;
  return r;
}
bool isprime(ull n) {
  if (n < 2 || n % 6 % 4 != 1) return (n | 1) == 3;
  ull A[] = {2, 325, 9375, 28178, 450775,
  9780504, 1795265022},
  s = __builtin_ctzll(n - 1), d = n >> s;
  for (ull a : A) {
    ull p = pow(a % n, d, n), i = s;
    while (p != 1 && p != n - 1 && a % n && i--)
      p = (ull)p * p % n;
    if (p != n - 1 && i != s) return 0;
  }
  return 1;
}
\end{minted}
\subsection{Pollard Rho (Fast Factorization)}
\begin{minted}{c++}
#define ull uint64_t
struct Montgomery {
  ull n, nr;
  Montgomery(ull n) {
    this->n = n, nr = 1;
    for (int i = 0; i < 6; i++)
      nr *= 2 - n * nr;
  }
  ull reduce(__uint128_t x) const {
    ull q = ull(x) * nr,
    m = ((__uint128_t)q * n) >> 64;
    return (x >> 64) + n - m;
  }
  ull multiply(ull x, ull y) {
    return reduce((__uint128_t)x * y); }
};
ull f(ull x, Montgomery& m, ull a) {
  return m.multiply(x, x) + a; }
ull diff(ull a, ull b) {
  return a > b ? a - b : b - a; }
const int M = 1024;
const ull SEED = 42;
ull find_factor(ull n, ull x0 = 2, ull a = 1) {
  Montgomery m(n);
  ull x = SEED;
  for (int l = M; l < (1 << 20); l *= 2) {
    ull y = x, p = 1;
    for (int i = 0; i < l; i += M) {
      for (int j = 0; j < M; j++)
        x = f(x, m, a), p = m.multiply(p, diff(x, y));
      if (ull g = gcd(p, n); g != 1) return g;
    }
  }
  return n;
}
\end{minted}
\subsection{Möbius Function and Inversion}
The Möbius function $\mu(n)$: \\
\textbullet{} $\mu(n) = 1$ if $n$ is square-free with even number of prime factors. \\
\textbullet{} $\mu(n) = -1$ if $n$ is square-free with odd number of prime factors. \\
\textbullet{} $\mu(n) = 0$ if $n$ has a squared prime factor.

\textbf{Properties \& Inversion}: \\
\textbullet{} $\sum_{d|n} \mu(d) = [n = 1]$ \\
\textbullet{} $g(n) = \sum_{d|n} f(d) \iff f(n) = \sum_{d|n} \mu(d)g(n/d)$ \\
\textbullet{} $\phi(n) = \sum_{d|n} \mu(d) \frac{n}{d}$ and $n = \sum_{d|n} d \cdot \mu\left(\frac{n}{d}\right)$ \\
\textbullet{} $\frac{n}{\phi(n)} = \sum_{d|n} \frac{\mu(d)}{d}$

\textbf{GCD Sums}: \\
\textbullet{} $\sum_{i=1}^n [\gcd(i, n) = k] = \phi\left(\frac{n}{k}\right)$ \\
\textbullet{} $\sum_{i=1}^n \gcd(i, n) = \sum_{d|n} d \cdot \phi\left(\frac{n}{d}\right)$ \\
\textbullet{} $\sum_{i=1}^n \frac{1}{\gcd(i,n)} = \frac{1}{n} \sum_{d|n} d \cdot \phi(d)$ \\
\textbullet{} $\sum_{i=1}^n \frac{i}{\gcd(i,n)} = \frac{1}{2} \sum_{d|n} d \cdot \phi(d)$

\subsection{GCD and LCM Properties}
\textbullet{} $\gcd(a, b) = \gcd(b, a \bmod b)$ \\
\textbullet{} If $a | (b \cdot c)$ and $\gcd(a, b) = d$, then $(a/d) | c$ \\
\textbullet{} $\gcd(a, \text{lcm}(b,c)) = \text{lcm}(\gcd(a,b), \gcd(a,c))$ \\
\textbullet{} $\gcd(n^a - 1, n^b - 1) = n^{\gcd(a,b)} - 1$

\subsection{Gauss Circle Theorem}
Lattice points inside $m^2 + n^2 \le r^2$: \\
$N(r) = 1 + 4 \sum_{i=0}^\infty \left( \lfloor \frac{r^2}{4i+1} \rfloor - \lfloor \frac{r^2}{4i+3} \rfloor \right)$
